\documentclass[aspectratio=169]{beamer}
\input{header}
\coursecode{JKSSB}

\title{PowerPoint Animation \& Transition}
\subtitle{Lesson 32 --- Effects and the Slideshow Keys}
\author{JKSSB Computer Awareness}
\date{\today}

\begin{document}
\frame{\titlepage}

\begin{frame}{Animation and Transition: The Core Difference}
  \begin{itemize}
    \item \key{Animation} is an effect applied to \key{one object or text}
      \key{inside a slide}.
    \item \key{Transition} is an effect applied to the movement
      \key{from one slide to the next}.
    \item Both are added from the \key{Animations} and \key{Transitions} tabs.
    \item They control \key{how things move}, not the content itself.
  \end{itemize}
  \begin{block}{The one idea}
    Animation works \key{within} a slide; a transition works \key{between}
    slides. Keep the two words apart.
  \end{block}
\end{frame}

\begin{frame}{What an Animation Applies To}
  \begin{itemize}
    \item An animation can be set on text, bullet points, pictures, shapes,
      charts, tables, and SmartArt.
    \item It can start \key{when the slide appears} or \key{on a mouse click}.
    \item Timing, duration, and order are controlled from the
      \key{Animation Pane}.
    \item Animation does not change the position of the slide in the deck;
      that is a transition's job.
  \end{itemize}
\end{frame}

\begin{frame}{The Four Animation Effect Categories}
  \begin{center}
    \begin{tabular}{p{0.20\textwidth}p{0.62\textwidth}}
      \toprule
      \textbf{Effect} & \textbf{What it does} \\
      \midrule
      Entrance & Makes an object \key{appear} on the slide. \\
      Exit & Makes an object \key{leave} the slide. \\
      Emphasis & Draws attention to an object that is \key{already visible}. \\
      Motion Path & Moves an object \key{along a defined path}. \\
      \bottomrule
    \end{tabular}
  \end{center}
  \vspace{0.2cm}
  \muted{Entrance and Exit change whether the object is on the slide;
  Emphasis changes how the object behaves; Motion Path changes where it
  travels.}
\end{frame}

\begin{frame}{Entrance versus Exit}
  \begin{columns}[T]
    \column{0.46\textwidth}
      \textbf{Entrance}
      \begin{itemize}
        \item Brings an object \key{onto} the slide.
        \item Used to reveal content step by step.
      \end{itemize}
    \column{0.46\textwidth}
      \textbf{Exit}
      \begin{itemize}
        \item Takes an object \key{off} the slide.
        \item Used to clear content before the next point.
      \end{itemize}
  \end{columns}
  \vspace{0.3cm}
  \begin{alertblock}{Near-neighbour alert}
    Entrance and Exit are opposites. An effect that makes an object stay and
    simply stand out is \key{Emphasis}, not Entrance.
  \end{alertblock}
\end{frame}

\begin{frame}{Emphasis and Motion Path}
  \begin{itemize}
    \item \key{Emphasis}: the object stays on the slide and is highlighted ---
      for example a grow, spin, or colour pulse.
    \item \key{Motion Path}: the object travels from one place to another
      along a drawn line, such as a straight line or a curve.
    \item A motion path changes an object's \key{position} during the show.
    \item Both belong to the \key{Animations} tab, not to Transitions.
  \end{itemize}
  \begin{exampleblock}{Quick check}
    ``The logo flies in from the left and settles into the corner''
    describes a \key{Motion Path}; ``the logo spins once and stops''
    describes \key{Emphasis}.
  \end{exampleblock}
\end{frame}

\begin{frame}{Transitions: Slide to Slide}
  \begin{itemize}
    \item A \key{transition} plays as the presentation moves from the
      \key{current slide to the next}.
    \item Standard transitions include \key{Fade}, \key{Push}, \key{Wipe},
      and \key{Dissolve}.
    \item A transition can be set with a \key{duration} and an optional
      \key{sound}.
    \item Transitions are managed on the \key{Transitions} tab and can be
      applied to one slide or to all slides.
  \end{itemize}
  \begin{block}{Applying to all slides}
    The \key{Apply To All} button copies the chosen transition to every slide
    in the presentation.
  \end{block}
\end{frame}

\begin{frame}{Transition Groups}
  \begin{center}
    \begin{tabular}{p{0.24\textwidth}p{0.58\textwidth}}
      \toprule
      \textbf{Group} & \textbf{Character} \\
      \midrule
      Subtle & Gentle, professional changes such as Fade and Wipe. \\
      Exciting & Stronger movement such as Push, Cover, and Split. \\
      Dynamic Content & Motion that responds to the slide content itself. \\
      \bottomrule
    \end{tabular}
  \end{center}
  \vspace{0.25cm}
  \muted{Transitions of the same group look consistent; mixing groups can make
  a deck look uneven.}
\end{frame}

\begin{frame}{The Morph Transition}
  \begin{itemize}
    \item \key{Morph} is a modern transition that \key{smoothly moves and
      animates} objects as one slide changes into the next.
    \item It works best when two slides share \key{similar objects} --- for
      example the same picture in two positions.
    \item Morph then makes those objects appear to \key{glide} into their new
      places.
    \item Morph is a \key{transition}, so it is set on the Transitions tab.
  \end{itemize}
  \begin{exampleblock}{Worked example}
    Slide 1 shows a circle on the left; slide 2 shows the same circle on the
    right. With \key{Morph}, the circle appears to slide smoothly across,
    instead of jumping.
  \end{exampleblock}
\end{frame}

\begin{frame}{Starting the Slide Show: F5}
  \begin{itemize}
    \item \key{F5} starts the slide show \key{from the beginning} --- from the
      first slide.
    \item \key{Shift+F5} starts the slide show \key{from the current slide}
      --- the one selected.
    \item \key{Alt+F5} starts the presentation in \key{Presenter View}.
    \item \key{Esc} \key{ends} a running slide show.
  \end{itemize}
  \begin{alertblock}{Real PYQ \textnormal{(WO 2026 Q80)}}
    ``Slide show from the beginning'' is \key{F5}. This exact pairing has
    appeared in the official paper.
  \end{alertblock}
\end{frame}

\begin{frame}{The Slideshow Key Set}
  \begin{center}
    \begin{tabular}{p{0.24\textwidth}p{0.58\textwidth}}
      \toprule
      \textbf{Key} & \textbf{Action during a slide show} \\
      \midrule
      F5 & Start from the \key{beginning}. \\
      Shift+F5 & Start from the \key{current} slide. \\
      Alt+F5 & Start in \key{Presenter View}. \\
      Esc & \key{End} the slide show. \\
      \bottomrule
    \end{tabular}
  \end{center}
  \vspace{0.25cm}
  \muted{Learn this table as a set; exam items often swap F5 and Shift+F5.}
\end{frame}

\begin{frame}{More Controls While Presenting}
  \begin{itemize}
    \item \key{N}, \key{Enter}, \key{Space}, or \key{Right Arrow}: go to the
      \key{next} slide or animation.
    \item \key{P} or \key{Left Arrow}: go back to the \key{previous} slide.
    \item \key{B}: turn the screen \key{black}; \key{W}: turn it \key{white}.
    \item Type a slide number and press \key{Enter} to jump to that slide.
  \end{itemize}
  \begin{block}{Presenter View}
    Presenter View shows the current slide, the \key{next} slide, the
    \key{speaker notes}, and a timer --- all on the presenter's screen.
  \end{block}
\end{frame}

\begin{frame}{Speaker Notes}
  \begin{itemize}
    \item The \key{Notes pane} sits below the slide in Normal view.
    \item It holds \key{speaker notes} --- reminders the presenter reads while
      presenting.
    \item Notes are \key{not shown} on the projected slide during a normal
      slide show.
    \item In \key{Presenter View}, the notes are visible only to the
      presenter.
  \end{itemize}
\end{frame}

\begin{frame}{PowerPoint Views}
  \begin{center}
    \begin{tabular}{p{0.26\textwidth}p{0.56\textwidth}}
      \toprule
      \textbf{View} & \textbf{Purpose} \\
      \midrule
      Normal & Edit one slide with an outline and the notes pane. \\
      Outline & Work with the text of all slides as an outline. \\
      Slide Sorter & See \key{thumbnails} of all slides and reorder them. \\
      Notes Page & View and print a slide with its speaker notes. \\
      Reading View & Review the show in a window without full screen. \\
      Slide Show & Present the deck full screen. \\
      \bottomrule
    \end{tabular}
  \end{center}
  \vspace{0.15cm}
  \muted{Slide Sorter is the view for rearranging slides; Notes Page is the
  view for notes with the slide.}
\end{frame}

\begin{frame}{Handouts}
  \begin{itemize}
    \item A \key{handout} prints the presentation for the \key{audience}.
    \item Handouts usually place \key{several slides on one page} --- two,
      three, four, six, or nine per page.
    \item Handout layouts can include \key{lined space} for the audience to
      write notes.
    \item Speaker notes are for the \key{presenter}; a handout is for the
      \key{audience}.
  \end{itemize}
  \begin{alertblock}{Do not confuse them}
    Notes pane = \key{presenter}. Handout = \key{audience}. They are different
    outputs.
  \end{alertblock}
\end{frame}

\begin{frame}{SmartArt Conversion and Its Limits}
  \begin{itemize}
    \item A \key{SmartArt} graphic can be converted to \key{shapes} or to
      \key{text}.
    \item \key{Convert to Text} turns the graphic into a \key{plain bulleted
      list}: the graphical layout is lost.
    \item \key{Convert to Shapes} keeps individual shapes but removes the
      SmartArt arrangement, so it cannot be turned back into SmartArt.
    \item Content that depended on the layout must be \key{re-entered} or
      re-applied by hand.
  \end{itemize}
  \begin{block}{Preview of a real PYQ}
    The JA 2026 paper tested which SmartArt content \key{cannot} convert
    cleanly without re-entry \textnormal{(PYQ-17)}. Remember: the \key{design
    and layout} are what a text conversion drops.
  \end{block}
\end{frame}

\begin{frame}{Trap Alert: F5 versus F7 versus F9}
  \begin{itemize}
    \item \key{F5} starts the slide show from the beginning.
    \item \key{Shift+F5} starts from the current slide.
    \item \key{Alt+F5} starts Presenter View.
    \item \key{F7} is the \key{spelling and grammar} check --- not the slide
      show (TRAP-32).
    \item \key{Esc} ends the show.
  \end{itemize}
  \begin{alertblock}{The trap}
    Do not choose F7 or F9 for ``start the slide show''. F5 is the show;
    F7 is spell check.
  \end{alertblock}
\end{frame}

\begin{frame}{Lesson 32: Recall}
  \begin{itemize}
    \item \key{Animation} acts on an object \key{within} a slide;
      \key{transition} acts \key{between} slides.
    \item Four animation effects: \key{Entrance}, \key{Exit},
      \key{Emphasis}, \key{Motion Path}.
    \item \key{Morph} is a transition that smoothly moves matching objects
      across slides.
    \item Show keys: \key{F5} = beginning, \key{Shift+F5} = current,
      \key{Alt+F5} = Presenter View, \key{Esc} = end.
    \item \key{Notes pane} = presenter; \key{handout} = audience;
      \key{Slide Sorter} = reorder slides.
    \item \key{SmartArt} converted to text loses its graphical layout.
  \end{itemize}
\end{frame}
\end{document}
